%%%%%%%% Ralphthon @ ICML 2026 — Track 1 AI Scientist (BLIND) %%%%%%%%
% 4-page hard limit (excl. references/appendix). Double-blind anonymized.
% All numbers trace to reports/logs/results_*.json (run 2026-07-12, seed 42, 200 personas/case).
\documentclass{article}
\usepackage{microtype}\usepackage{graphicx}\usepackage{booktabs}\usepackage{hyperref}
\newcommand{\theHalgorithm}{\arabic{algorithm}}
\usepackage{icml2026}          % blind; camera-ready: \usepackage[accepted]{icml2026}
\usepackage{amsmath}\usepackage{amssymb}\usepackage{mathtools}
\usepackage[capitalize,noabbrev]{cleveref}
\icmltitlerunning{Do Synthetic Personas Predict Real Conversion?}
\begin{document}
\twocolumn[
  \icmltitle{Do Synthetic Personas Predict Real Conversion?\\
  A Blind, Leakage-Controlled Backtest Across Three Decision Framings}
  \icmlsetsymbol{equal}{*}
  \begin{icmlauthorlist}\icmlauthor{Anonymous Author(s)}{anon}\end{icmlauthorlist}
  \icmlaffiliation{anon}{Anonymous Institution}
  \icmlcorrespondingauthor{Anonymous}{anon@example.com}
  \icmlkeywords{synthetic personas, demand validation, agent simulation, backtest, LLM agents}
  \vskip 0.3in
]
\printAffiliationsAndNotice{}

\begin{abstract}
Synthetic personas are increasingly proposed as a cheap substitute for real
demand validation, yet their ability to predict \emph{real-world} conversion is
seldom tested against ground truth. We build a blind, leakage-controlled
backtest---five documented historical validation episodes, anonymized so an LLM
cannot recall their outcomes---and have persona agents from a population-grounded
dataset make purchase decisions, scoring against the real result. Across 200
personas per case we find naive one-shot persona simulation is
\emph{anti-predictive}: real successes receive \emph{lower} simulated conversion
than real failures (success$-$failure separation $\Delta=-0.46$). Crucially, this
failure is \textbf{framing-invariant}: reformulating the decision from stated
willingness to a longitudinal revealed-preference judgment collapses every
response to a floor ($\Delta=0$), and a graded-probability formulation restores
resolution but not discrimination ($\Delta=-0.05$); the case ranking is identical
across all three. The failure is also \textbf{aggregation-invariant}: real
successes are dominated by failures not only in the mean but at every quantile
including the top decile and the maximum, so no hidden enthusiast minority explains
the real winners. Applying the method to two \emph{real} pre-launch products
reproduces the bias (predicting $\sim$14\% paid conversion for a product whose
team targets 5\%). Yet the bias is partly \emph{correctable}: leave-one-out
calibration on real outcomes raises directional accuracy from $0.40$ (raw) to
$0.64$ over 14 cases, though a perfect result on the initial five did not survive
expansion---the sim--real map is \emph{non-monotonic}, since some products
(over-hyped hardware) are correctly rejected by the personas. We conclude
synthetic-persona sims must not be used raw and only partly recover under simple
calibration, and we release the blind backtest as a reusable validity gate.
\end{abstract}

\section{Introduction}
\label{sec:intro}
AI coding lets founders ship faster than ever, yet few products sustain value;
demand validation (interviews, fake-door tests, A/B) is the recommended hedge but
is costly for solo builders. Synthetic personas promise to automate it, and
recent work offers population-grounded persona datasets~\citep{nemotron} and
multi-agent simulation engines~\citep{mirofish}. The unanswered question is
epistemic: \emph{can a persona simulation predict conversion that holds in the
real world?} We contribute (i) a \textbf{blind, leakage-controlled backtest
protocol} isolating predictive signal from LLM memorization via case
anonymization; (ii) a \textbf{three-framing robustness study} (stated, revealed,
graded) showing the failure is not an artifact of one prompt; and (iii) a
\textbf{cautionary real-product application} to two live pre-launch products; and
(iv) a \textbf{calibration analysis} showing the bias is partly correctable
(leave-one-out accuracy rises from $0.40$ to $0.64$ over 14 cases) but that the
sim--real map is non-monotonic, so simple calibration is insufficient. The raw
finding is negative, obtained without tuning toward any desired outcome.

\section{Related Work}
\label{sec:related}
\emph{Silicon sampling}~\citep{argyle2023} conditions an LLM on a respondent's
backstory to reproduce survey distributions, and now underlies most ``AI panel''
products. Its validity is uneven: reviews of silicon-to-human comparisons find
only a minority match while most diverge, and willingness-to-pay studies report
directional alignment in \emph{familiar} categories but weak performance in
\emph{novel} ones~\citep{persona_microdata,persona_reliability}. The documented
weak spots---novel-category behavior, minority-opinion tails, and actual purchase
in unfamiliar contexts---are precisely where demand validation operates.
Population-grounded datasets~\citep{nemotron} and multi-agent
engines~\citep{mirofish} scale personas, but prior work largely evaluates
\emph{plausibility} or survey-distribution match. We instead test \emph{external
predictive validity} against real conversion outcomes under a blind,
leakage-controlled protocol, and quantify how far light calibration closes the
gap---corroborating and sharpening the WTP-in-novel-categories weakness.

\section{Method}
\label{sec:method}
\paragraph{Personas.} We sample personas from a $\sim$1M-persona
population-grounded Korean dataset~\citep{nemotron} (occupation, education, region,
consumption narrative), rendered to a compact profile. A local 7B
instruction-tuned LLM~\citep{qwen25} is prompted to \emph{become} the persona.
\paragraph{Three decision framings.} \textbf{Stated}: one-shot ``would you pay?''
(binary). \textbf{Revealed}: a longitudinal judgment---click, try, and still pay
after 2--8 weeks once novelty fades and incumbents are considered (binary).
\textbf{Graded}: a continuous probability of becoming a retained paying customer.
A case's simulated conversion is the mean over its personas.
\paragraph{Persona conditions.} \textbf{General}: random personas.
\textbf{Target}: personas keyword-filtered to each product's target segment, fixed
\emph{a priori} from the product description (never the outcome).
\paragraph{Baseline and leakage control.} A no-persona baseline asks the LLM for
the conversion rate directly. Each case is anonymized (brand, company, era
removed) so the model cannot recall the historical result; ground truth is
revealed only for scoring. We report directional accuracy and separation
$\Delta=\overline{\text{sim}}_{\text{success}}-\overline{\text{sim}}_{\text{failure}}$
($\Delta>0$ predictive, $\Delta<0$ anti-predictive).

\section{Results}
\label{sec:results}
\Cref{tab:cases} gives per-case simulated conversion under the graded framing;
real successes (A,B,E) sit \emph{below} real failures (C,D). \Cref{tab:framings}
shows this is framing-invariant: no framing or persona condition yields $\Delta>0$.

\begin{table}[t]\caption{Graded per-case conversion (200 personas, general).
Successes A,B,E score at or below failures C,D. Ranking C$>$D$>$B$>$A$>$E is
identical under all three framings.}\label{tab:cases}\vskip 0.1in
\centering\small\begin{tabular}{llcc}\toprule
Case & Method & GT & Graded \\ \midrule
A & fake-door video/waitlist & 1 & 0.153 \\
B & concierge fake store & 1 & 0.202 \\
C & reformulate flagship & 0 & 0.230 \\
D & Wizard-of-Oz & 0 & 0.180 \\
E & landing-page pricing & 1 & 0.126 \\ \bottomrule
\end{tabular}\vskip 0.1in\end{table}

\begin{table}[t]\caption{Separation $\Delta$ across three framings and two persona
conditions. Every value is $\le 0$; no configuration is predictive. Revealed
collapses to a floor (all responses $0$).}\label{tab:framings}\vskip 0.1in
\centering\small\begin{tabular}{lcc}\toprule
Framing & General $\Delta$ & Target $\Delta$ \\ \midrule
No-persona baseline & \multicolumn{2}{c}{$-0.233$} \\
Stated (binary) & $-0.462$ & $-0.494$ \\
Revealed (binary) & $0.000$ & $0.000$ \\
Graded (continuous) & $-0.045$ & $-0.054$ \\ \bottomrule
\end{tabular}\vskip 0.1in\end{table}

Directional accuracy never exceeds the no-persona baseline ($0.40$). The
\emph{stated} framing over-predicts appealing-but-failed cases; the \emph{revealed}
framing is so strict that every persona declines (a degenerate floor); the
\emph{graded} framing restores resolution but leaves failures ranked above
successes. The invariance of the ranking across framings, persona conditions, and
against the baseline indicates a structural, not prompt-specific, failure.

\paragraph{Aggregation robustness.} One might suspect the mean hides an intense
niche minority for the real successes. It does not: recomputing per-case
distributions, the failures dominate the successes on the high-intensity fraction,
the top-decile mean ($\Delta=-0.21$), the standard deviation, and the maximum
($\Delta=-0.32$). No moment of the persona distribution is predictive, closing the
``wrong statistic'' objection.

\paragraph{Real-product application.} We apply the graded method to two live
pre-launch products (anonymized): \textbf{P1}, a real-time civic-debate app with a
paid subscription tier, and \textbf{SoloSquad}~\citep{solosquad}, an open-source
multi-agent tool for solo founders. Predicted paid conversion is $0.14$ (P1) and
$0.10$ (SoloSquad), essentially
unchanged under target conditioning. P1's own team targets a $5\%$ premium
conversion; the simulation over-predicts by $\sim$3$\times$, reproducing the
abstract-appeal bias on real inputs.

\paragraph{Calibration partly recovers prediction.} The bias is systematic: real
successes sit consistently below real failures ($\Delta=-0.18$ over 14 cases). We
exploit this with leave-one-out calibration---a one-feature decision stump fit on
held-out cases. On the initial five cases LOO reaches $5/5$, but expanding to 14
(adding well-documented successes and failures) drops LOO to $0.64$---still above
raw ($0.40$) and chance ($0.50$), yet far from perfect. The residual errors are
informative: over-hyped expensive hardware is \emph{correctly} rejected by the
personas (low sim, real failure), breaking the anti-correlation and making the
sim--real map non-monotonic. Simple sign-flip calibration therefore helps but does
not solve the problem; adding an a-priori price/hype regime feature does not
improve LOO (also $0.64$), so the ceiling appears set by data scale, not feature
choice---reliable validity would require a substantially larger labeled corpus.
Pre-launch products without real labels (P1, SoloSquad) still require a real anchor
(e.g., a fake-door test) to calibrate.

\section{Discussion: Why It Fails}
\label{sec:discuss}
\textbf{(1) Abstract-appeal bias.} One-shot willingness elicits stated appeal, not
revealed conversion; broadly-appealing-but-failed products (a sweeter drink,
effortless dictation) score high. \textbf{(2) Missing temporal/relational
dynamics.} The real failures hinge on brand attachment and sustained-use
fatigue---invisible to a single decision; forcing a longitudinal binary instead
floors everything. \textbf{(3) Counterintuitive niche demand.} Real successes had
low general appeal but strong niche demand, which target filtering and graded
probabilities still underrate. Because the same ranking survives every
reformulation, the limitation is intrinsic to persona willingness, not to a
particular prompt. The bias is systematic enough that
calibration against real outcomes (\cref{sec:results}) \emph{partly} recovers
prediction, but a non-monotonic sim--real map---some duds are correctly rejected---
limits any simple correction: the practical recommendation is to treat persona
simulations as a \emph{feature} requiring regime-aware calibration, never a raw
estimate.

\section{Limitations}
\label{sec:limits}
Five backtest cases yield low statistical power; we frame results as a
mechanism-focused, framing-robust case study, not a definitive effect size. We use
one 7B model, one (Korean) persona market, and keyword target definitions; larger
models, multi-turn deliberation, and revealed-behavior traces are untested and are
exactly what our diagnosis motivates. The real-product predictions have no ground
truth and are reported only to illustrate the bias.

\section*{Impact Statement}
This work cautions against substituting synthetic-persona simulations for real
user evidence in product decisions; over-trust could bias which products get built
and entrench dataset demographics. By reporting a rigorous, framing-robust negative
result we aim to temper deployment claims. No human subjects or private data are
involved; product references are anonymized.

\section*{Software and Data}
The persona-simulation harness, prompts, anonymized case corpus, and per-persona
logs for all framings are released as part of the open-source SoloSquad
project~\citep{solosquad} at \url{https://github.com/Adelie-Squad/solosquad}
(logs under \texttt{reports/logs/}). Author identity is withheld for review.

\bibliography{references}
\bibliographystyle{icml2026}

\newpage\appendix\onecolumn
\section{Anonymized Case Corpus}
A device-syncing tool via demo-video waitlist (real: success); B online shoe store
via concierge fake store (success); C reformulated flagship drink that won blind
taste tests, replacing the original for loyal customers (failure); D voice
dictation assessed by Wizard-of-Oz for sustained use (low demand); E social-post
scheduler via landing-page pricing (success). Ground truth is revealed only for
scoring.
\section{Framings, Prompts, and Targets}
Stated elicits \texttt{would\_pay}; revealed elicits \texttt{retained\_pay} after a
click$\rightarrow$try$\rightarrow$retain funnel with incumbent attachment; graded
elicits \texttt{pay\_prob}$\in[0,1]$. Targets were fixed a priori by keyword (e.g.,
E: marketing/creator/self-employed terms) matched against occupation and
persona-narrative fields. Matched target counts (8000-pool): A 3062, B 711, C 2034,
D 2834, E 1023. Full prompts and keyword lists are in the repository.
\section{Per-Framing Logs}
Per-persona decisions and free-text reasons are in
\texttt{results\_\{stated,revealed,graded\}\_\{general,target\}.json} and, for real
products, \texttt{results\_products\_graded\_*.json}.
\end{document}
